% discussion.tex
\section{Discussion}
\label{sec:discussion}

\paragraph{The detectors carry little leakage information that loss does not.}
\citet{hayes2025strong} already report that membership-attack success and extraction come apart. Our
contribution is the measurement-validity question that follows: the contamination signal which
predicts \emph{extraction} is, to the resolution of our experiment, \emph{just raw loss}. The reference-free detectors that the contamination-detection
literature has invested in, Min-K\%, Min-K\%++, zlib, improve membership ranking by re-calibrating
the per-token likelihood (z-scoring against the vocabulary, compressing, or trimming to the
lowest-probability tokens), and one reading consistent with our data is that these adjustments
remove loss-magnitude information that tracks extractability. We cannot separate that reading from
suppression under near-collinearity, so we do not assert it as a mechanism. A descriptive mediation decomposition is consistent with
this, the loss-mediated (indirect) path is positive for all three detectors while the direct paths
are null or negative, but we read it descriptively, not causally: the detectors are near-collinear
transforms of loss (Spearman up to $0.90$, VIF up to $6.2$), so a negative direct/partial term is
consistent with statistical suppression rather than genuine inverse prediction. We therefore claim
only the conservative version: the calibrated detectors add \emph{no positive} leakage signal beyond
loss. A practitioner who wants to know \emph{which contaminated items the model will actually leak}
is, on this evidence, no better served by a state-of-the-art membership detector than by raw loss.
This gives a mechanistic form to the divergence that \citet{hayes2025strong} first reported, on the
reference-free detectors an auditor can actually afford.

\paragraph{Why this is a security result, not a leaderboard result.} Our finding is deliberately
\emph{not} ``we built a better detector.'' It is that the privacy question, will contamination of a
benchmark expose a leakage channel? is mis-served by importing the membership-inference toolkit
wholesale. For an auditor of a released model, the actionable implication is to measure
loss/extractability directly and to treat a high Min-K\%/Min-K\%++ score as evidence about
membership, not about leakage risk. The contribution is the assumption chain of
Section~\ref{sec:chain} plus this controlled measurement of incremental value. The detectors
themselves are prior work. Every statement in this section is bounded by the regime it was measured
in: one $160$M model, $N=300$ members, three fully extracted items, and membership AUC at chance. It
is a bound on what these instruments demonstrably carry, not a demonstration that they carry
nothing at scale.

\paragraph{Relation to prior work.} Our direction agrees with published results and we do not
claim the bottom line is surprising. \citet{hayes2025strong} established the divergence between
membership success and extraction, and \citet{alsahili2025effectiveness} report only ``marginal''
gains of MIA scores over likelihood ranking for targeted extraction. What we add is the controlled
form of the claim, a pre-registered partial-correlation/mediation
that quantifies a \emph{zero-to-negative} residual for the calibrated reference-free detectors after
loss is removed, and we target the reference-free detectors the contamination literature actually
deploys rather than a shadow-model attack. \citet{chen2025statistical} independently find
these detectors do not robustly beat the loss baseline for \emph{membership} once seed variance is
accounted for. Our result is the extraction-outcome analogue.

\paragraph{Defenses.} Because the leakage we measure is downstream of memorization, the principled
mitigation is differential privacy applied at training time~\citep{abadi2016deep,li2022dpllm}. It is a
producer-side control, not an auditor-side detector, and bounds the very quantity (loss-magnitude /
memorization) our analysis identifies as the operative one.
